\documentclass[11pt,oneside,openany]{book}
\usepackage[T1]{fontenc}
\usepackage[utf8]{inputenc}
\usepackage[sc]{mathpazo}
\renewcommand{\sfdefault}{lmss}
\renewcommand{\ttdefault}{lmtt}
\usepackage{amsmath,amssymb,amsthm,mathtools}
\usepackage[paperwidth=7.5in,paperheight=10in,inner=.85in,outer=.85in,top=.85in,bottom=.85in,headheight=15pt,headsep=20pt]{geometry}
\usepackage{microtype}
\usepackage{xcolor,graphicx,booktabs,array,longtable}
\newcolumntype{L}[1]{>{\raggedright\arraybackslash}p{#1}}
\usepackage{tikz,pgfplots}
\usetikzlibrary{arrows.meta,calc,patterns,decorations.markings}
\pgfplotsset{compat=1.18}
\definecolor{ink}{HTML}{17334B}
\definecolor{teal}{HTML}{137C83}
\definecolor{copper}{HTML}{AE5932}
\definecolor{soft}{HTML}{EEF3F5}
\pgfplotsset{every axis/.append style={axis line style={black!65},tick style={black!60},tick label style={font=\small},label style={font=\small},legend style={font=\small,draw=none,fill=white},grid style={black!10},major grid style={black!10},line width=.6pt},every axis plot/.append style={line width=1.15pt}}
\usepackage{enumitem,needspace,etoolbox}
\setlist[enumerate]{topsep=5pt,itemsep=5pt,parsep=0pt}
\setlist[itemize]{topsep=4pt,itemsep=3pt,parsep=0pt}
\usepackage{titlesec,fancyhdr}
\titleformat{\chapter}[display]{\normalfont\color{ink}}{\large\sffamily\MakeUppercase{Chapter \thechapter}}{12pt}{\huge\bfseries\raggedright\hyphenpenalty=10000\exhyphenpenalty=10000}[\vspace{12pt}\titlerule]
\titlespacing*{\chapter}{0pt}{-10pt}{24pt}
\titleformat{\section}{\large\bfseries\color{ink}}{\thesection}{.65em}{}
\titleformat{\subsection}{\normalsize\bfseries\color{ink}}{\thesubsection}{.6em}{}
\titlespacing*{\section}{0pt}{18pt plus 3pt minus 2pt}{7pt}
\titlespacing*{\subsection}{0pt}{12pt plus 2pt minus 2pt}{5pt}
\pagestyle{fancy}
\fancyhf{}
\fancyhead[L]{\small\nouppercase{\leftmark}}
\fancyfoot[C]{\small\thepage}
\renewcommand{\headrulewidth}{.3pt}
\renewcommand{\chaptermark}[1]{\markboth{\thechapter\quad #1}{}}
\fancypagestyle{plain}{\fancyhf{}\fancyfoot[C]{\small\thepage}\renewcommand{\headrulewidth}{0pt}}
\usepackage[font=small,labelfont=bf,labelsep=period]{caption}
\usepackage{float,placeins}
\floatplacement{figure}{htbp}
\setcounter{topnumber}{3}
\setcounter{bottomnumber}{2}
\renewcommand{\topfraction}{.9}
\renewcommand{\textfraction}{.08}
\renewcommand{\floatpagefraction}{.8}
\usepackage{makeidx}
\makeindex
\usepackage[pdfencoding=auto,psdextra,colorlinks=true,linkcolor=ink,citecolor=teal,urlcolor=teal,bookmarksnumbered=true]{hyperref}
\usepackage{bookmark}
\hypersetup{pdftitle={Discrete Information Geometry},pdfsubject={Probability, counting, and finite dynamics}}
\urlstyle{same}
\newtheorem{theorem}{Theorem}[chapter]
\newtheorem{proposition}[theorem]{Proposition}
\newtheorem{lemma}[theorem]{Lemma}
\newtheorem{corollary}[theorem]{Corollary}
\theoremstyle{definition}
\newtheorem{definition}[theorem]{Definition}
\newtheorem{example}{Example}[chapter]
\BeforeBeginEnvironment{example}{\par\addvspace{4pt}\noindent\begin{minipage}{\linewidth}}
\AfterEndEnvironment{example}{\end{minipage}\par\addvspace{4pt}}
\theoremstyle{remark}
\newtheorem*{remark}{Remark}
\newenvironment{solution}{\par\smallskip\noindent\textit{Solution.}\ }{\par\medskip}
\newenvironment{exercises}{\begin{enumerate}[label=\thechapter.\arabic*,ref=\thechapter.\arabic*,leftmargin=2.8em,itemsep=6pt]}{\end{enumerate}}
\newenvironment{answers}{\begin{description}[leftmargin=2.9em,labelwidth=2.5em,labelsep=.4em,itemsep=8pt,font=\normalfont\bfseries]}{\end{description}}
\newcommand{\R}{\mathbb{R}}
\newcommand{\E}{\mathbb{E}}
\newcommand{\Var}{\operatorname{Var}}
\newcommand{\Cov}{\operatorname{Cov}}
\newcommand{\KL}[2]{D\!\left(#1\,\middle\|\,#2\right)}
\newcommand{\grad}{\operatorname{grad}}
\newcommand{\arcosh}{\operatorname{arcosh}}
\newcommand{\sech}{\operatorname{sech}}
\newcommand{\softmax}{\operatorname{softmax}}
\newcommand{\diag}{\operatorname{diag}}
\newcommand{\dd}{\,\mathrm{d}}
\newcommand{\ip}[2]{\left\langle #1,#2\right\rangle}
\newcommand{\p}{\mathbf{p}}
\newcommand{\q}{\mathbf{q}}
\newcommand{\f}{\mathbf{f}}
\newcommand{\uvec}{\mathbf{u}}
\newcommand{\vvec}{\mathbf{v}}
\newcommand{\ones}{\mathbf{1}}
\newcommand{\simp}{\Delta^{n-1}}
\newcommand{\simpi}{\Delta^{n-1}_{\!\circ}}
\numberwithin{equation}{chapter}
\setlength{\parindent}{1.2em}
\setlength{\parskip}{0pt}
\linespread{1.07}
\setlength{\emergencystretch}{2em}
\clubpenalty=10000
\widowpenalty=10000
\displaywidowpenalty=10000
\raggedbottom
\setcounter{tocdepth}{1}
\setcounter{secnumdepth}{2}
\providecommand{\tightlist}{\setlength{\itemsep}{0pt}\setlength{\parskip}{0pt}}
\providecommand{\passthrough}[1]{#1}
\newcommand{\labheading}{\FloatBarrier\section*{Laboratory investigation}\addcontentsline{toc}{section}{Laboratory investigation}}

\begin{document}
\begin{titlepage}\color{ink}\centering\vspace*{1.0in}{\small\sffamily PROBABILITY AND MATHEMATICAL MODELLING\par}\vspace{.45in}{\Huge\bfseries Discrete Information\\[.12in]Geometry\par}\vspace{.35in}{\Large Probability, Counting,\\and Finite Dynamics\par}\vspace{1cm}{\Large Chapter 3\par}\vfill\begin{tikzpicture}[scale=1.05]\draw[teal,line width=1pt](0,0)--(4,0)--(2,3.464)--cycle;\foreach \a in {0,...,6}{\foreach \b in {0,...,6}{\pgfmathtruncatemacro{\c}{6-\a-\b}\ifnum\c>-1\fill[teal] ({(\b+\c/2)*4/6},{\c*3.464/6}) circle(1.5pt);\fi}}\end{tikzpicture}\vfill{\large A course in ten chapters\par}\vspace{.3in}{\small\url{https://bobovski66.github.io/IntroDiscreteInfoGeom/}\par}\end{titlepage}
\setcounter{chapter}{2}

\chapter{Three strategies and probability triangles}
Three types allow motion in two independent directions. A triangle provides a complete picture of their proportions, including mixtures on its edges and pure states at its vertices. Integer population counts supply a natural lattice inside the triangle. These two representations connect geometry with combinatorics.

\hypertarget{ch03-reading-a-probability-triangle}{%
\section{Reading a probability triangle}\label{ch03-reading-a-probability-triangle}}

Write \(p=(p_A,p_B,p_C)\) with nonnegative entries adding to one. Place the pure states at \(A=(0,0)\), \(B=(1,0)\), and \(C=(1/2,\sqrt3/2)\). Represent \(p\) by the weighted point \[X=p_AA+p_BB+p_CC=\left(p_B+\frac{p_C}{2},\frac{\sqrt3}{2}p_C\right).\] These are \textbf{barycentric coordinates}: the coefficients tell us how much of each vertex contributes. The midpoint of edge \(AB\) is \((1/2,1/2,0)\) in probability coordinates. The center is \((1/3,1/3,1/3)\). A point on edge \(AB\) has \(p_C=0\).

To recover probabilities from a plotted point \((u,v)\), solve the two coordinate equations: \[p_C=\frac{2v}{\sqrt3},\qquad p_B=u-\frac{v}{\sqrt3},\qquad p_A=1-u-\frac{v}{\sqrt3}.\] The point is inside the triangle precisely when these three values are nonnegative.

\begin{example}[Locating a state]

The distribution \((1/2,1/3,1/6)\) is plotted at \((5/12,\sqrt3/12)\). Its height is one sixth of the triangle's altitude, reflecting its \(C\) proportion. Its position is closer to vertex \(A\) than to either of the other vertices, consistent with \(A\) being the most common type.

\end{example}

\hypertarget{ch03-tokens-produce-a-triangular-lattice}{%
\section{Tokens produce a triangular lattice}\label{ch03-tokens-produce-a-triangular-lattice}}

Suppose the total count is a fixed integer \(N\). A count triple \((a,b,c)\) with \(a+b+c=N\) gives the probability vector \((a/N,b/N,c/N)\). The possible states form a lattice. Moving one token from \(A\) to \(B\) changes the counts by \((-1,1,0)\) and the probabilities by \((-1/N,1/N,0)\).

There are \[\frac{(N+1)(N+2)}2=\binom{N+2}{2}\] possible triples. One proof fixes \(c\): for \(c=0,1,\ldots,N\), there are \(N-c+1\) choices for \(a\), with \(b\) then determined. Adding gives \((N+1)+N+\cdots+1\). A second proof encodes the triple as \(N\) stars separated into three groups by two bars; choosing the two bar positions among \(N+2\) positions determines the triple.

\begin{example}[A lattice of fifteen states]

With \(N=4\), the lattice has \(\binom62=15\) states. The count triple \((2,1,1)\) is adjacent to six states obtained by moving a token between two distinct types. At \((4,0,0)\), only two moves are possible because an empty type cannot donate a token.

\end{example}

A normalized growth update generally does not land on this lattice. It defines a deterministic map on all probability vectors. Sampling \(N\) new individuals from the updated proportions produces a random lattice state instead. These are two distinct models with different sources of variation.

\begin{figure}[htbp]\centering
\resizebox{.98\linewidth}{!}{\begin{tikzpicture}\draw[ink,thick] (0,0)--(6,0)--(3,5.196)--cycle;\node[below left] at (0,0){$A$};\node[below right] at (6,0){$B$};\node[above] at (3,5.196){$C$};\fill[teal] (3,5.19615) circle (2pt);\fill[teal] (3.75,3.89711) circle (2pt);\fill[teal] (4.5,2.59808) circle (2pt);\fill[teal] (5.25,1.29904) circle (2pt);\fill[teal] (6,0) circle (2pt);\fill[teal] (2.25,3.89711) circle (2pt);\fill[teal] (3,2.59808) circle (2pt);\fill[teal] (3.75,1.29904) circle (2pt);\fill[teal] (4.5,0) circle (2pt);\fill[teal] (1.5,2.59808) circle (2pt);\fill[teal] (2.25,1.29904) circle (2pt);\fill[teal] (3,0) circle (2pt);\fill[teal] (0.75,1.29904) circle (2pt);\fill[teal] (1.5,0) circle (2pt);\fill[teal] (0,0) circle (2pt);\node[anchor=west,align=left] at (6.8,3.7){$N=4$: fifteen states\\One token changes two counts.};\draw[copper,thick] (2.25,1.29904) circle (5pt);\node[below=6pt] at (2.25,1.29904){$(2,1,1)$};\draw[copper,->] (2.25,1.29904)--(3.75,1.29904);\draw[copper,->] (2.25,1.29904)--(3,2.59808);\draw[copper,->] (2.25,1.29904)--(0.75,1.29904);\draw[copper,->] (2.25,1.29904)--(1.5,2.59808);\draw[copper,->] (2.25,1.29904)--(1.5,0);\draw[copper,->] (2.25,1.29904)--(3,0);\end{tikzpicture}}
\caption{The denominator-four token lattice. Arrows show the six one-token moves from counts (2,1,1).}\end{figure}

\hypertarget{ch03-cyclic-interactions}{%
\section{Cyclic interactions}\label{ch03-cyclic-interactions}}

For three strategies, a score table is a \(3\) by \(3\) matrix. The average score of type \(i\) is \(f_i=\sum_jM_{ij}p_j\). Consider \[M_s=\begin{pmatrix}0&1&-s\\-s&0&1\\1&-s&0\end{pmatrix},\qquad s>0.\] Type \(A\) gains against \(B\), type \(B\) gains against \(C\), and type \(C\) gains against \(A\). Each suffers a loss of magnitude \(s\) against the remaining opponent. At the center, all three scores equal \((1-s)/3\), so the center is fixed for both normalized multiplier rules.

For exponential multipliers, a step is \[p_i^+=\frac{p_i e^{\eta f_i}}{Z},\qquad Z=\sum_jp_j e^{\eta f_j}.\] The current scores must be recomputed before the next step. Using the same scores repeatedly instead would describe fixed external advantages, rather than this interaction game.

\begin{example}[Computing a cyclic step]

Let \(s=1\) and \(p=(1/2,1/3,1/6)\). The scores are \((1/6,-1/3,1/6)\). With \(\eta=1\), cancelling the common factor \(e^{1/6}\) leaves unnormalized entries \((1/2,e^{-1/2}/3,1/6)\). After normalization the result is approximately \((0.5755,0.2327,0.1918)\). The two equally scoring types retain their ratio \(3:1\) during this step.

\end{example}

\hypertarget{ch03-a-finite-inequality-detects-outward-drift}{%
\section{A finite inequality detects outward drift}\label{ch03-a-finite-inequality-detects-outward-drift}}

The \textbf{weighted arithmetic-geometric mean inequality} states that, for positive \(a_i\) and nonnegative weights \(p_i\) adding to one, \[\sum_i p_i a_i\ \geq\ \prod_i a_i^{p_i}.\] For rational weights, it follows from the ordinary arithmetic-geometric mean inequality by repeating each entry the appropriate number of times. The real-weight version follows by approximation. Equality holds when all entries carrying positive weight are equal. We will use this finite inequality several times.

For the balanced cyclic table \(s=1\), direct cancellation gives \[\sum_i f_i=0,\qquad \sum_i p_if_i=0.\] Let \(P(p)=p_Ap_Bp_C\). Under exponential normalization, \[P(p^+)=\frac{P(p)e^{\eta(f_A+f_B+f_C)}}{Z^3}=\frac{P(p)}{Z^3}.\] Applying weighted arithmetic-geometric mean to \(a_i=e^{\eta f_i}\) gives \(Z\geq e^{\eta\sum_i p_if_i}=1\). For an interior state other than the center, the scores are not all equal, so \(Z>1\) and \(P\) strictly decreases.

The product is largest at the center: the ordinary arithmetic-geometric mean inequality gives \(p_Ap_Bp_C\leq(1/3)^3\). Decreasing product therefore records movement through product contours toward the boundary. It does not imply that every individual coordinate moves monotonically, nor does it by itself identify a limit point. A trajectory may continue circulating while this diagnostic decreases.

\begin{proposition}[The balanced finite rule does not conserve the product]

For every positive intensity, an exponential update of the balanced cyclic game strictly decreases \(p_Ap_Bp_C\) at every interior state other than the center.

\end{proposition}

There is a similar conclusion for linear multipliers whenever they are positive. Since the weighted mean score is zero, their normalizing denominator is \(1\). Their unweighted sum is \(3\), so their product is at most \(1\). Multiplying the three updated coordinates shows \(P(p^+)\leq P(p)\), again strictly away from the center. These exact calculations explain why changing the size of a finite step can affect a cyclic picture.

\hypertarget{ch03-comparing-a-different-recurrence}{%
\section{Comparing a different recurrence}\label{ch03-comparing-a-different-recurrence}}

The permutation rule \[T(p_A,p_B,p_C)=(p_C,p_A,p_B)\] also moves states cyclically. It satisfies \(T^3(p)=p\) and preserves the coordinate product exactly. Its trajectories contain at most three states. This recurrence is useful as a simple example of exact periodicity and as a check on a plotting program. It is a different update from score-based selection, with a different algebraic structure.

\begin{figure}[htbp]\centering
\resizebox{.98\linewidth}{!}{\begin{tikzpicture}\draw[ink,thick] (0,0)--(6,0)--(3,5.196)--cycle;\node[below left] at (0,0){$A$};\node[below right] at (6,0){$B$};\node[above] at (3,5.196){$C$};\draw[teal,line width=1pt,] plot coordinates {(2.5,0.866025) (2.30211,0.914991) (2.13433,0.989415) (2.00434,1.09113) (1.91704,1.22169) (1.87454,1.3817) (1.87635,1.57002) (1.91959,1.78302) (1.99934,2.014) (2.10912,2.25316) (2.2416,2.48824) (2.38957,2.70569) (2.54686,2.89222) (2.70911,3.03607) (2.87405,3.12795) (3.04129,3.16149) (3.21161,3.13339) (3.38597,3.04357) (3.56442,2.8955) (3.74508,2.69646) (3.92355,2.45754) (4.09295,2.19293) (4.24443,1.91836) (4.3683,1.64902) (4.45503,1.39764) (4.4961,1.17313) (4.48445,0.98029) (4.41464,0.820371) (4.28327,0.692054) (4.08965,0.592573) (3.83693,0.518653) (3.53324,0.467199) (3.19222,0.435747) (2.83231,0.422736) (2.4746,0.427675) (2.13975,0.451239) (1.84514,0.495337) (1.60317,0.563116) (1.42101,0.65886) (1.30151,0.787681) (1.24443,0.954869) (1.2474,1.16479) (1.30623,1.41932) (1.41474,1.71605) (1.56431,2.04684) (1.744,2.39749) (1.94142,2.74913) (2.14444,3.08112) (2.34321,3.37439) (2.5317,3.61395) (2.70818,3.78973) (2.87485,3.89591) (3.03679,3.92947) (3.20062,3.88899) (3.37309,3.7742) (3.55961,3.58665) (3.76275,3.33128) (3.98096,3.01809) (4.20804,2.66311) (4.4337,2.28732) (4.64535,1.91367) (4.83045,1.56301) (4.97842,1.25062) (5.08144,0.984702) (5.13404,0.766899) (5.13197,0.594087) (5.07111,0.460494) (4.94694,0.359436) (4.75496,0.28447) (4.49214,0.229994) (4.15941,0.191484) (3.76431,0.165501) (3.32271,0.149618) (2.85811,0.142316) (2.39816,0.142908) (1.96937,0.151513) (1.59225,0.169073) (1.27893,0.197432) (1.03359,0.239463) (0.85476,0.299217) (0.738192,0.382049) (0.679138,0.494619) (0.673725,0.644618) (0.719298,0.839994) (0.813842,1.08747) (0.954673,1.39033) (1.13681,1.74578) (1.35164,2.14306) (1.58679,2.56335) (1.82756,2.98255) (2.05966,3.37612) (2.27206,3.72397) (2.4587,4.01334) (2.61864,4.23892) (2.75509,4.40077) (2.87403,4.50162) (2.98296,4.54436) (3.0901,4.53041) (3.20382,4.45901) (3.33225,4.3274) (3.48251,4.13207) (3.65971,3.87086) (3.86549,3.5458) (4.09666,3.16584) (4.34467,2.74806) (4.5966,2.31618) (4.83773,1.89618) (5.05474,1.51085) (5.23811,1.17555) (5.38285,0.896891) (5.48765,0.673976) (5.55317,0.501037) (5.58036,0.370035) (5.56908,0.272596) (5.51723,0.201121) (5.42044,0.149249) (5.27216,0.111934) (5.0646,0.0853168) (4.79047,0.0665215) (4.44576,0.053457) (4.0331,0.044639) (3.56458,0.0390514) (3.06221,0.0360444) (2.555,0.0352672) (2.07288,0.0366313) (1.64035,0.0403013) (1.27234,0.0467149) (0.973804,0.0566307) (0.74195,0.0712098) (0.569614,0.0921336) (0.448243,0.121762) (0.369969,0.163332) (0.328758,0.221175) (0.320881,0.300924) (0.344943,0.409613) (0.401587,0.555512) (0.492862,0.74747) (0.621167,0.993495) (0.78772,1.29846) (0.990736,1.6612) (1.22388,2.07208) (1.47591,2.51243) (1.73205,2.95717) (1.9771,3.38005) (2.19865,3.75943) (2.38931,4.08194) (2.54698,4.34304) (2.67379,4.54499) (2.7745,4.6939) (2.855,4.79702) (2.92146,4.86085) (2.97988,4.89007) (3.03604,4.88695) (3.09566,4.85123) (3.16458,4.78006) (3.24888,4.66828) (3.35478,4.50906) (3.48814,4.29509) (3.65339,4.02072) (3.85203,3.68474) (4.08103,3.29339)};
\begin{axis}[at={(7cm,0.2cm)},anchor=south west,width=6.2cm,height=5.2cm,xlabel={Step $t$},ylabel={$p_Ap_Bp_C$},xmin=0,xmax=160,ymin=0]
\addplot[teal] coordinates {(0,0.0277778) (1,0.0275013) (2,0.0272369) (3,0.0269819) (4,0.0267289) (5,0.026468) (6,0.0261894) (7,0.0258869) (8,0.0255602) (9,0.025215) (10,0.0248622) (11,0.0245144) (12,0.024182) (13,0.02387) (14,0.0235766) (15,0.0232934) (16,0.0230078) (17,0.0227064) (18,0.0223788) (19,0.0220213) (20,0.0216382) (21,0.021241) (22,0.0208458) (23,0.0204676) (24,0.0201169) (25,0.019796) (26,0.0194993) (27,0.0192144) (28,0.0189254) (29,0.0186171) (30,0.0182787) (31,0.0179075) (32,0.01751) (33,0.0171006) (34,0.0166977) (35,0.0163182) (36,0.0159732) (37,0.0156654) (38,0.0153888) (39,0.0151309) (40,0.0148755) (41,0.014606) (42,0.0143097) (43,0.0139805) (44,0.0136212) (45,0.0132435) (46,0.0128648) (47,0.0125037) (48,0.0121745) (49,0.0118844) (50,0.0116325) (51,0.0114107) (52,0.0112064) (53,0.0110044) (54,0.0107901) (55,0.0105519) (56,0.0102839) (57,0.00998763) (58,0.00967252) (59,0.00935355) (60,0.00904742) (61,0.00876811) (62,0.00852376) (63,0.00831568) (64,0.00813921) (65,0.00798544) (66,0.00784306) (67,0.00770007) (68,0.00754534) (69,0.00737022) (70,0.00717027) (71,0.00694668) (72,0.00670663) (73,0.00646191) (74,0.00622596) (75,0.00601042) (76,0.00582259) (77,0.00566453) (78,0.00553376) (79,0.00542456) (80,0.00532941) (81,0.00524001) (82,0.00514812) (83,0.00504626) (84,0.0049286) (85,0.00479199) (86,0.00463703) (87,0.00446848) (88,0.00429461) (89,0.00412528) (90,0.00396945) (91,0.00383319) (92,0.00371888) (93,0.00362559) (94,0.00355002) (95,0.00348759) (96,0.00343313) (97,0.0033814) (98,0.00332729) (99,0.00326618) (100,0.00319423) (101,0.00310906) (102,0.00301042) (103,0.00290072) (104,0.00278492) (105,0.00266959) (106,0.00256129) (107,0.00246505) (108,0.00238349) (109,0.00231688) (110,0.0022637) (111,0.00222139) (112,0.00218692) (113,0.00215719) (114,0.00212915) (115,0.00209989) (116,0.00206665) (117,0.00202701) (118,0.00197908) (119,0.00192197) (120,0.00185626) (121,0.00178419) (122,0.00170947) (123,0.00163651) (124,0.00156934) (125,0.00151076) (126,0.00146195) (127,0.0014227) (128,0.00139184) (129,0.00136771) (130,0.0013485) (131,0.00133244) (132,0.00131786) (133,0.00130316) (134,0.00128681) (135,0.00126735) (136,0.00124346) (137,0.0012141) (138,0.00117883) (139,0.00113805) (140,0.0010932) (141,0.00104661) (142,0.00100108) (143,0.000959129) (144,0.000922544) (145,0.0008921) (146,0.000867697) (147,0.000848643) (148,0.000833956) (149,0.000822582) (150,0.000813507) (151,0.000805804) (152,0.00079862) (153,0.000791151) (154,0.000782608) (155,0.000772206) (156,0.000759189) (157,0.000742912) (158,0.000722983) (159,0.000699445) (160,0.000672936)};
\end{axis}\end{tikzpicture}}
\caption{Balanced cyclic exponential updates at intensity 0.35. The triangle trajectory circulates while the coordinate product decreases.}\end{figure}

\hypertarget{ch03-reading-distances-and-contours-in-a-triangle}{%
\section{Reading distances and contours in a triangle}\label{ch03-reading-distances-and-contours-in-a-triangle}}

Barycentric coordinates have another useful interpretation. The perpendicular distance from a point to side \(AB\), divided by the triangle's full altitude, is \(p_C\). To verify this, note that the side is horizontal and that the point's height is \((\sqrt3/2)p_C\). By symmetry, the same statement holds for the other two coordinates and their opposite sides. Lines of constant \(p_C\) are therefore parallel to side \(AB\).

\begin{example}[Recovering a composition from height ratios]

Suppose a point lies at one fifth of the full altitude from side \(AB\) and at one half of the full altitude from side \(BC\). The first measurement gives \(p_C=1/5\), and the second gives \(p_A=1/2\). The remaining coordinate is \(p_B=3/10\). The three normalized distances add to one because they are barycentric coordinates.

\end{example}

A contour of a function is a set of states with the same function value. For the coordinate product, a contour is described by \(p_Ap_Bp_C=c\). The product is zero on the boundary and at most \(1/27\) in the interior. Product contours surround the center, but a one-step decrease in product does not specify a unique direction along a contour or a unique next coordinate. A scalar diagnostic compresses a two-dimensional state into one number, so it necessarily discards some detail.

\hypertarget{ch03-counting-the-interior-and-the-boundary}{%
\section{Counting the interior and the boundary}\label{ch03-counting-the-interior-and-the-boundary}}

The stars-and-bars argument extends to \(k\) types. Nonnegative integer counts summing to \(N\) are encoded by \(N\) stars and \(k-1\) bars, so their number is \(\binom{N+k-1}{k-1}\). Empty groups correspond to zero counts. If every count must be positive, first give one token to each type. The remaining \(N-k\) tokens can be distributed freely, giving \(\binom{N-1}{k-1}\) states when \(N\geq k\).

For three types, the number of boundary states is the total minus the strictly positive count. For \(N\geq3\), this is \[\binom{N+2}{2}-\binom{N-1}{2}=3N.\] One can also count \(N+1\) points on each of three edges and subtract three to correct the double counting of vertices. The two answers agree. Comparing two counting arguments is a useful way to verify a combinatorial formula.

\begin{example}[A larger lattice]

At \(N=8\), there are \(\binom{10}{2}=45\) total states and \(\binom72=21\) strictly positive states. The remaining \(24\) are on the boundary. A uniform distribution over these 45 lattice vertices would be a probability model on compositions. It is different from independently sampling eight individuals from a uniform three-type population: the latter assigns each count vector a multinomial weight.

\end{example}

\hypertarget{ch03-sampling-can-change-support}{%
\section{Sampling can change support}\label{ch03-sampling-can-change-support}}

At a positive state, deterministic positive multiplication keeps every type present. A sample of \(N\) individuals can omit a type by chance. On an edge of the triangle, a sampling model without mutation cannot restore the missing type, because its sampling probability is zero. A finite sampled population may therefore reach an edge and subsequently evolve on that edge, even when a deterministic trajectory remains in the interior for every finite step.

This observation explains why a deterministic trajectory and a simulated finite population should be labeled separately. Both can be useful descriptions. The deterministic rule isolates state-dependent reweighting; the sampled model also includes fluctuations associated with finitely many individuals. Chapter 4 provides exact probabilities for those fluctuations.

\hypertarget{ch03-exercises}{%
\section{Exercises}\label{ch03-exercises}}

\begin{exercises}

\item 

Plot \((1,0,0)\), \((0,1/2,1/2)\), and \((1/4,1/4,1/2)\) using barycentric coordinates.

\item 

Recover the probabilities of the point \((u,v)=(1/2,\sqrt3/6)\).

\item 

List all count triples with \(N=2\). Verify the formula for their number.

\item 

How many lattice states are there for \(N=10\)? How many have all three counts positive?

\item 

List the neighbors of \((2,1,1)\) under one-token moves.

\item 

Compute the scores of the balanced cyclic table at \((1/5,1/2,3/10)\).

\item 

Verify algebraically that \(\sum_i p_if_i=0\) for the balanced table.

\item 

Derive the product update formula for exponential multipliers for any three scores, without assuming balance.

\item 

Explain why the product test says nothing new on an edge of the triangle. Determine whether the update can leave that edge.

\item 

Show that the center is the unique interior state where all balanced cyclic scores agree.

\item 

For the permutation rule, determine all fixed states and all possible periods.

\item 

Compare deterministic normalization with sampling a fixed number of individuals after normalization. Which can lose a type in a single step when every initial proportion is positive?

\end{exercises}

\hypertarget{ch03-selected-solutions}{%
\section{Selected solutions}\label{ch03-selected-solutions}}

\textbf{2.} The recovered probabilities are \((1/3,1/3,1/3)\), the center.

\textbf{4.} There are \(\binom{12}{2}=66\) states. Subtract one token from each positive count; the remaining three nonnegative counts sum to \(7\), giving \(\binom92=36\) strictly positive states.

\textbf{5.} The neighbors are \((1,2,1)\), \((1,1,2)\), \((3,0,1)\), \((2,0,2)\), \((3,1,0)\), and \((2,2,0)\).

\textbf{6.} The scores are \((1/5,1/10,-3/10)\).

\textbf{8.} Multiplying the three update equations gives \(P(p^+)=P(p)e^{\eta\sum_i f_i}/Z^3\). Balance is needed only to simplify the exponential factor to one.

\textbf{11.} A fixed state has three equal coordinates and is therefore the center. Every other state has period three, since applying the permutation three times is the identity and a period dividing three is either one or three.
\labheading\begin{enumerate}
\def\labelenumi{\arabic{enumi}.}
\tightlist
\item
  Plot the states \((1,0,0)\), \((1/2,1/2,0)\), \((1/3,1/3,1/3)\), and \((1/2,1/3,1/6)\). Calculate their Cartesian coordinates and label the absent types on each edge.
\item
  Draw the denominator-four token lattice. Count its states in two ways and list all neighbors of \((2,1,1)\). Distinguish a one-token move from a general normalized growth step.
\item
  For the balanced cyclic table, derive the three scores at \((1/2,1/3,1/6)\). Verify that their sum and their probability-weighted mean are both zero. Calculate one exponential step with intensity \(1\).
\item
  Run exponential cyclic updates at intensities \(0.1\), \(0.5\), and \(1\) for fifty steps from the same state. Record the initial and final coordinate products and inspect the triangle traces. Use the exact product identity to explain the direction of change.
\item
  Compare with the permutation rule. List its first four states and prove its period and product preservation. Explain why the permutation calculation does not establish product preservation for score-based selection.
\end{enumerate}

\FloatBarrier

\end{document}
